\textit{Reality is simultaneously preserved within and performed by the body; this establishes the body as an archive of cultural knowledge and memory. One that can also serve as a driving force as well.
}

\subsection{Black Performance Epistemologies}(\parencite{Patrick_2003}, \parencite{DeFrantz_2006}, \parencite{DeFrantz_Gonzalez_2014})
Performance has often served as an epistemological practice that generates cultural knowledge and understanding through embodied experience. This has been no different for Black and other marginalized communities. The analysis of Black movement has led to understanding of how these movements encode and transmit cultural memory. Even providing a framework for how these performances crate alternative ways of knowing being and doing.

Bridge:

These embodied ways of knowing find parallel expression in Black sonic traditions, which often function as their own sophisticated technological systems.

\subsection{Sonic Traditions as Technical Practice}(\parencite{Gilroy_1993}, \parencite{Weheliye_2005}, \parencite{Weheliye_2014}, \parencite{weheliye2005})
As previously mentioned, musical and sonic practices are foundational to Black diasporic cultures; these traditions function as technologies used to transmit culture across geographic and temporal boundaries. These technologies are key to the temporal alternative and strengthen social relationships. They can also function as biopolitical technologies that have been used as both constraints and resistance for Black life and experience.

Bridge:

While sonic practices demonstrate one form of embodied technological knowledge, broader research in multisensory computing reveals how these principles can inform digital design methodologies overall

\subsection{Embodied Interaction and Multisensory Computing}(\parencite{Hook_2018}, \parencite{Jones_Arning_Farver_2006}, \parencite{Sha_2013})
Work in somaesthetic interaction design provides some foundational understanding of how embodied knowledge informs technological designs. There is a long history of practitioners who have explored are refuted dominant sensory hierarchies through art and technology that can provide a lens for better understanding embodied and multisensory experiences. And these practices and processes indeed do create new technological possibilities through their explorations of space and interactions.

Bridge:

Embodied knowledge systems require community structures for their preservation and transmission, pointing toward collective models that influence both  individual and group technological creation.

\subsection{Community Knowledge and Collective Creation}(\parencite{Smith_2012}, \parencite{Escobar_2018})
There exists work in trying to better navigate multiple cultural system simultaneously that can inform how to retain authenticity within dominant technological systems. This also requires that the research and implementation processes be maintained through community control. This can lead to alternative forms of technologies rooted in the marginalized communities relational practices, without a need for unnecessary extraction.

Bridge to Next Section:

These community-centered approaches to knowledge creation become particularly crucial when examining how cultural assumptions become embedded within computational systems themselves.